% Generated draft: requires mathematical and source review.
\documentclass[11pt]{article}
\usepackage[margin=25mm]{geometry}
\usepackage{fontspec}
\setmainfont{Latin Modern Roman}
\usepackage{amsmath,amssymb,mathtools,unicode-math}
\setmathfont{Latin Modern Math}
\usepackage{enumitem,needspace,xurl,hyperref}
\setlength{\parindent}{0pt}
\setlength{\parskip}{5pt}
\setlength{\emergencystretch}{3em}
\newcommand{\N}{\mathbb{N}}
\newcommand{\Z}{\mathbb{Z}}
\newcommand{\Q}{\mathbb{Q}}
\newcommand{\R}{\mathbb{R}}
\newcommand{\C}{\mathbb{C}}
\newcommand{\F}{\mathbb{F}}
\newcommand{\A}{\mathbb{A}}
\newcommand{\PP}{\mathbb{P}}
\newcommand{\E}{\mathbb{E}}
\newcommand{\eps}{\varepsilon}
\newcommand{\dd}{\,\mathrm{d}}
\begin{document}
\section*{UnsolvedMath ID: 20000381}
\textbf{Source problem code:} catalogue entry
\subsection{Definitions and mathematical statement}
\textbf{Original question (Jordan Ellenberg).} Understand which ways of counting field extensions satisfy Principles (A) and (B). Neither principle is defined or assigned a mathematical meaning in the supplied source record. Thus a precise classification under those exact labels cannot be supplied without additional context.

\textbf{Attached unreviewed reconstruction and criteria.} The aid cautiously proposes two context-supported targets, a one-place Chebotarev law (Ch) and finite-place asymptotic independence (Ind), without assigning either to (A) or (B). For a family of extensions with nonnegative weights and height cutoff $T$, let $N(T)$ be total weight and $N_S(T)$ the weight satisfying chosen local events $E_v$ at a fixed finite place set $S$. Assume a common scale $\Phi(T)$ and fixed-local-set error control giving
\[N(T)=C_0\Phi(T)+o(\Phi(T)),\qquad N_S(T)=C_S\Phi(T)+o(\Phi(T)),\qquad C_0>0.\]
Write $C_v$ for a single-event constant. The reported criterion is
\[\mathrm{Ind}\quad\Longleftrightarrow\quad C_SC_0^{|S|-1}=\prod_{v\in S}C_v\quad\text{for every fixed finite local-event pattern}.\]
This is asymptotic independence, not exact independence at finite cutoff. The corresponding interaction defect is $C_SC_0^{|S|-1}-\prod_vC_v$; its nonvanishing detects failure. The division-free version accommodates zero-density individual events.

For a reconstructed unramified Chebotarev target with fixed group $G$, conjugacy class $c\subseteq G$, and one-place constants $C_{v,c}$ and $C_{v,\mathrm{unr}}$, the expected conditional class ratio is $|c|/|G|$. The constant test is
\[|G|C_{v,c}=|c|C_{v,\mathrm{unr}}.\]
A conditional probability additionally needs $C_{v,\mathrm{unr}}>0$; if that constant is zero, the equality alone is vacuous. The source does not specify the extension family, group, ramification conditioning or height, so this formula describes the aid's diagnostic framework rather than an identification of the missing principles.

Under a separate common Tauberian theorem carrying dominant Laurent coefficients to counting constants by the same nonzero factor for all fixed local conditions, the aid says the same homogeneous tests can be applied to those Laurent coefficients. Such a transfer requires the analytic hypotheses and common dominant singularity; none is automatic from a formal Euler product.

The aid also reports two global failure mechanisms. For a coherent persistent two-component mixture of product local laws, with limiting weight $0<\theta<1$ and component local probabilities $u_v,w_v$, the covariance for two local indicators is
\[\theta(1-\theta)(u_v-w_v)(u_{v'}-w_{v'}).\]
A nonzero value obstructs Ind. For a finite product probability law $\mu=\bigotimes_v\mu_v$ tilted by a nonnegative integrable density $W$ with $0<Z=\int W\,\dd\mu<\infty$, the product-tilt criterion is that the tilted law remains a product precisely when $W/Z$ factors $\mu$-almost everywhere into normalized one-coordinate densities. Zero-probability conditioning is excluded; dependence in a finite-level count need not persist asymptotically.

The aid says these abstract criteria were compared with Wood's abelian fair-count theorem, discriminant counterexamples, and a cyclic-prime full-genus theorem, but gives no complete statements or references. It expressly does not establish the difficult family-specific asymptotics needed to classify counting orders generally. All criteria and comparisons require independent review.
\subsection{Short English statement}
Ellenberg asks which extension-counting orders satisfy two undefined principles. The aid reconstructs Chebotarev and independence targets without identifying the letters, then gives conditional leading-constant tests and mixture/conditioning obstructions. These diagnostics do not supply the missing definitions or prove the required counting asymptotics for arbitrary families.
\subsection{Sources}
\begin{itemize}
\item Local UnsolvedMath registry, permanent ID 20000381, source code catalogue entry; \texttt{aim-algebraic-number-theory-notes.json}, notes[32], tag \texttt{question}.
\item AIM workshop \emph{Arithmetic statistics over finite fields and function fields}, January 2014, Question 3 (Jordan Ellenberg): \url{https://aimath.org/pastworkshops/arithstatffieldproblems.pdf}. Principles (A) and (B) are not defined in the supplied record. No PDF was consulted.
\item Attached research attempt 1 names Wood's abelian fair-count theorem, discriminant counterexamples and a cyclic-prime full-genus theorem without bibliographic locators. Result \texttt{partial\_result}, status \texttt{partially\_solved}, rigor label \texttt{proved}, novelty confidence low. It is a machine-generated, non-peer-reviewed aid requiring independent verification before citation. Original literature field empty.
\end{itemize}
\end{document}
